\documentclass[11pt]{article}
\usepackage[margin=1in]{geometry}
\usepackage{amsmath,amssymb,amsthm}
\usepackage[colorlinks=true,urlcolor=blue]{hyperref}
\newtheorem{theorem}{Theorem}
\newtheorem{lemma}{Lemma}
\newtheorem{proposition}{Proposition}
\newcommand{\Fp}{\mathbb F_p}
\newcommand{\Prob}{\mathbb P}
\newcommand{\Ent}{\mathsf H}
\newcommand{\KL}{\mathsf D}
\title{Arbitrary-target configuration counts and quantitative entropy tails}
\author{Research working note: direct adaptation, historical priority unverified}
\date{November 15, 2026}
\begin{document}
\maketitle

Fix integers $d\ge3$ and an odd prime $p$. The results below are uniform in
every target vector but keep $d,p$ fixed. Let $C$ be the adjacency matrix
of the uniform pairing model on $N$ vertices with $d$ labelled stubs per
vertex; assume only $dN$ even. Loops contribute two to the diagonal.
In particular, odd $N$ is allowed when $d$ is even.

Huang's published-topic preprint \cite{huang}, Theorem 1.1(1.2), concerns
the zero-target kernel moment, not the all-target mixing statement needed
for the odd-block RSBM extension. This note proves the arbitrary-target
enumeration and entropy-tail parts directly. The remaining central local
limit theorem is not asserted here. The combinatorial enumeration follows
the mechanism of Huang's Proposition 2.2, with different row targets;
the maximization over targets and the resulting rate function are proved
below rather than imported as a theorem from that source.

\section{Exact prescribed-output enumeration}
For $q\in\Fp^N$, let
\[
 N_a=|\{v:q_v=a\}|,\qquad \nu_a=N_a/N \quad(a\in\Fp).
\]
If a pairing has adjacency matrix $C$, its directed color-count matrix is
\[
 M_{ab}=\sum_{v:q_v=a}\sum_{w:q_w=b} C_{vw}.
\]
Thus $M$ is symmetric, every $M_{aa}$ is even, and
$\sum_b M_{ab}=dN_a$. Put $m=M/(dN)$, so $m$ is a symmetric probability
matrix with both marginals $\nu$. Define
\[
 P_t(z)=\sum_{\substack{a_1,\ldots,a_d\in\Fp\\a_1+\cdots+a_d=t}}
                \prod_{j=1}^d z_{a_j},
\qquad T_{dN}=(dN-1)!!.
\]
The coefficient notation $[z^{M_a}]$ uses $p$ variables and the row
$M_a=(M_{ab})_b$.

\begin{proposition}[Exact arbitrary-output formula]\label{prop:exact}
For every fixed $q,r\in\Fp^N$ and admissible matrix $M$,
\begin{align}
 &\Prob(Cq=r,\ M(q,C)=M)\\
 &\quad=\frac{K(M)}{T_{dN}}
       \prod_{a\in\Fp}[z^{M_a}]
              \prod_{v:q_v=a}P_{r_v}(z),\label{eq:count}\\
 &K(M)=\prod_{a<b} M_{ab}!\,
         \prod_a\frac{M_{aa}!}{2^{M_{aa}/2}(M_{aa}/2)!}.
 \nonumber
\end{align}
\end{proposition}
\begin{proof}
For each stub at vertex $v$, record the color of its partner's vertex.
At $v$, the ordered $d$ partner colors have sum $r_v$ if and only if
$(Cq)_v=r_v$. The coefficient in \eqref{eq:count} counts assignments of
these colors to the labelled stubs with specified total row counts.
For a fixed assignment with symmetric counts $M$, the stubs from color
$a$ requesting color $b$ can be paired bijectively with those from $b$
requesting $a$ in $M_{ab}!$ ways. The $M_{aa}$ stubs requesting their own
color have $(M_{aa}-1)!!$ pairings. Multiplication gives $K(M)$.
Every pairing has exactly one such assignment, and the total number of
pairings is $T_{dN}$. This proves the formula, including loops and
parallel edges.
\end{proof}

\section{Uniform row inequality and equality cases}
Set $\alpha=(d-1)/d$. For a probability vector $\pi$ on $\Fp$, define
\[
 F_t(\pi)=\sum_{a_1+\cdots+a_d=t}\prod_{j=1}^d\pi(a_j)^\alpha,
 \qquad F_*(\pi)=\max_{t\in\Fp}F_t(\pi).
\]
Let $u=(1/p,\ldots,1/p)$.

\begin{lemma}\label{lem:row}
For every probability vector $\pi$,
\[
 p^{-(d-1)}\le F_*(\pi)\le1.
\]
Moreover $F_t(\pi)=1$ if and only if either $\pi=u$, or $\pi$ is
a point mass at $a$ and $t=da$ in $\Fp$. No assumption $p\nmid d$
is needed for this lemma.
\end{lemma}
\begin{proof}
For every tuple, the arithmetic--geometric mean inequality gives
\[
 \prod_{j=1}^d\pi(a_j)^{(d-1)/d}
 \le\frac1d\sum_{j=1}^d\prod_{\ell\ne j}\pi(a_\ell).
\]
After summing over tuples with prescribed sum $t$, each term obtained
by omitting one coordinate sums to one. Thus $F_t\le1$.
The tuple with every entry equal to a largest atom $a$ gives
$F_{da}(\pi)\ge\pi(a)^{d-1}\ge p^{-(d-1)}$.

For equality, let $S=\operatorname{supp}\pi$. If $d-1$ entries lie in
$S$ and the completing entry lies outside $S$, the corresponding AM--GM
inequality is strict. Hence equality requires
\[
 t-(d-1)S\subseteq S.
\]
Fixing $d-2$ entries in $S$ gives a reflection preserving $S$.
Changing one of these fixed entries from $a$ to $b$ and composing the
two reflections gives invariance of $S$ under translation by $a-b$.
If $|S|\ge2$, a nonzero such translation generates the additive group
of $\Fp$, so $S=\Fp$. Since $d\ge3$, any two values can appear together
in a tuple of prescribed sum. AM--GM equality for that tuple forces
their probabilities to agree. Thus $\pi=u$. For $|S|=1$, direct
substitution gives the stated point-mass case. The two asserted cases
also plainly give equality.
\end{proof}

\begin{lemma}[Uniform Taylor expansion]\label{lem:taylor}
On the hyperplane $\sum_a z_a=0$, as $z\to0$,
\[
 \log F_t(u+z)
 =-\frac{(d-1)p}{2d}\|z\|_2^2+O_{d,p}(\|z\|_2^3),\label{eq:taylor}
\]
uniformly over $t\in\Fp$. Consequently, on a sufficiently small fixed
neighborhood of $u$,
\[
 \log F_*(\pi)\le-\frac{(d-1)p}{4d}\|\pi-u\|_2^2.
\]
\end{lemma}
\begin{proof}
Let $(A_1,\ldots,A_d)$ be uniform on tuples of sum $t$. Since $d\ge3$,
each coordinate is uniform on $\Fp$ and every pair is independent,
for every $t$. Then
\[
 F_t(u+z)=\mathbb E\prod_{j=1}^d(1+pz_{A_j})^\alpha.
\]
The linear term vanishes. Pairwise independence makes every quadratic
cross term vanish. The remaining quadratic coefficient is
\[
 \frac{\alpha(\alpha-1)}2p^2d\,
       \mathbb E z_{A_1}^2
 =-\frac{(d-1)p}{2d}\sum_a z_a^2.
\]
On $\max_a|z_a|\le1/(2p)$ all these functions have bounded third
derivatives. There are only $p$ targets, so the remainder is uniform.
The logarithm has the same quadratic term. Shrinking the neighborhood
absorbs the cubic remainder into half the negative quadratic term.
\end{proof}

\section{Rate bound, with every polynomial factor charged}
For a symmetric probability matrix $m$ with row marginal $\nu$, define
$\pi_a=m_a/\nu_a$ when $\nu_a>0$. Define
\[
 \Psi(m)=-\frac{d-2}{2}\KL(m\|\nu\otimes\nu)
            +\sum_{a:\nu_a>0}\nu_a\log F_*(\pi_a).
 \tag{3}\label{eq:rate}
\]
Rows of zero mass contribute zero. The relative entropy is finite:
if $m_{ab}>0$, then $\nu_a\nu_b>0$. As usual $0\log0=0$.

\begin{proposition}\label{prop:typebound}
There are constants $C=C(d,p)$ and an integer
$K_0=4p^2+20p+20$ (a deliberately loose choice) such that, for every
admissible type $M$, every input $q$ of profile $\nu$, and every target $r$,
\begin{align}
 \Prob(Cq=r,M(q,C)=M)
 &\le C(N+1)^{K_0}
          \exp\{N[-\Ent(\nu)+\Psi(m)]\},\label{eq:fixedq}\\
 \sum_{q:\operatorname{type}(q)=\nu}
       \sup_r\Prob(Cq=r,M(q,C)=M)
 &\le C(N+1)^{K_0}\exp\{N\Psi(m)\}.
 \label{eq:summedq}
\end{align}
The prefactors are uniform over the targets; no enumeration of target
types is performed.
\end{proposition}
\begin{proof}
Evaluate each coefficient in \eqref{eq:count} at
$z_b=\pi_a(b)^\alpha$ to obtain
\[
 [z^{M_a}]\prod_{v:q_v=a}P_{r_v}(z)
 \le F_*(\pi_a)^{N_a}\prod_b\pi_a(b)^{-\alpha M_{ab}}.
 \tag{4}\label{eq:coeff}
\]
If a coordinate of $\pi_a$ is zero, its $M_{ab}$ is also zero;
the bound follows by taking positive-coordinate limits. Empty rows
give coefficient one.

Uniform factorial bounds give
\[
 \log\frac{K(M)}{T_{dN}}
   =\frac{dN}{2}\sum_{a,b}m_{ab}\log m_{ab}
           +O_{d,p}(\log(N+1)).
\]
Indeed apply $\log j!=j\log j-j+O(\log(dN+1))$ to each factorial,
including $j=0$. There are only $O(p^2)$ factorials. Adding the
logarithm of \eqref{eq:coeff} gives the exponent
\[
 \frac d2\sum m_{ab}\log m_{ab}
  -(d-1)\sum m_{ab}\log\frac{m_{ab}}{\nu_a}
  +\sum_a\nu_a\log F_*(\pi_a)
  =-\Ent(\nu)+\Psi(m).
\]
This proves \eqref{eq:fixedq}. There are
$N!/\prod_aN_a!=\exp\{N\Ent(\nu)+O_p(\log(N+1))\}$ inputs of
that profile. Multiplying proves \eqref{eq:summedq}.
The displayed loose exponent $K_0$ dominates these factorial errors;
all $d$-dependent constant factors are absorbed into $C(d,p)$.
\end{proof}

\begin{lemma}[Zeros and quantitative uniformity]\label{lem:rate}
The function $\Psi$ is continuous and nonpositive on the compact set
of symmetric probability matrices. Its zeros are exactly
\[
 m=u\otimes u\quad\hbox{and}\quad m=e_a\otimes e_a\quad(a\in\Fp).
\]
For each fixed $\eta>0$ there exists $c=c(d,p,\eta)>0$ such that
\[
 \max_a\nu_a\le1-\eta
 \quad\Longrightarrow\quad
 \Psi(m)\le-c\bigl(\|m-u\otimes u\|_2^2+\|\nu-u\|_2^2\bigr).
 \tag{5}\label{eq:quadratic}
\]
If the domain is empty, the assertion is vacuous.
\end{lemma}
\begin{proof}
The bound $p^{-(d-1)}\le F_*\le1$ shows that each zero-row weighted
term in \eqref{eq:rate} tends to zero continuously. Also
\[
 \KL(m\|\nu\otimes\nu)=\sum m_{ab}\log m_{ab}-2\sum\nu_a\log\nu_a
\]
is continuous. Nonpositivity follows from Lemma~\ref{lem:row}.
If $\Psi=0$, positivity of $d-2$ forces $m=\nu\otimes\nu$.
Every nonempty conditional row then equals $\nu$, and
Lemma~\ref{lem:row} forces $\nu=u$ or a point mass. The converse
is immediate.

In a sufficiently small neighborhood of $u\otimes u$, all marginals
are bounded below, and all conditional rows lie in the neighborhood
from Lemma~\ref{lem:taylor}. Put
\[
 Q=\sum_a\nu_a\|\pi_a-u\|_2^2.
\]
Symmetry implies $\sum_a\nu_a\pi_a=\nu$, so Jensen gives
$Q\ge\|\nu-u\|_2^2$. Moreover,
\[
 \|m-\nu\otimes u\|_2^2
   =\sum_a\nu_a^2\|\pi_a-u\|_2^2\le Q.
\]
The rows of $m-\nu\otimes u$ sum to zero, so they are orthogonal
to the corresponding constant rows of $(\nu-u)\otimes u$. Thus
\[
 \|m-u\otimes u\|_2^2
 =\|m-\nu\otimes u\|_2^2+p^{-1}\|\nu-u\|_2^2
 \le(1+p^{-1})Q.
\]
Consequently the sum of the two squared distances in
\eqref{eq:quadratic} is at most $3Q$, while
\[
 \Psi(m)\le-\frac{(d-1)p}{4d}Q.
\]
This proves a uniform local quadratic bound. On the remaining
compact portion with $\max\nu_a\le1-\eta$, none of the zeros of
$\Psi$ remains. The negative maximum there, divided by an upper
bound on the squared distances, extends the inequality to the
whole required domain.
\end{proof}

\begin{theorem}[All-target noncentral tails]\label{thm:tails}
Fix $\eta>0$ and $A>0$. There exists a finite constant
$B=B(d,p,\eta,A)$ such that
\begin{multline}
 \sum_{\substack{q\in\Fp^N\\\max_a\nu_a(q)\le1-\eta}}
 \sup_{r\in\Fp^N}\Prob\biggl(Cq=r,\\
 \|m(q,C)-u\otimes u\|_2^2+\|\nu(q)-u\|_2^2
       >B\frac{\log(N+1)}N\biggr)
 =O_{d,p,\eta,A}(N^{-A}).\label{eq:tailtheorem}
\end{multline}
In particular this gives an arbitrary-target summed maximum-atom
bound outside the shrinking balanced-input window, once
nearconstant inputs are excluded.
\end{theorem}
\begin{proof}
For each $q$, move the target supremum inside the sum over types,
which only enlarges the expression. Apply
\eqref{eq:summedq} and \eqref{eq:quadratic}. There are at most
$(N+1)^p$ input profiles and $(dN+1)^{p(p+1)/2}$ symmetric count
matrices before imposing any row or parity restrictions. Hence the
left side is bounded by
\[
 C_{d,p}(N+1)^{K_1-cB},\qquad
 K_1=K_0+p+p(p+1)/2.
\]
Choose $B>(K_1+A+1)/c$. The argument used only the condition
$dN$ even, and all constants are independent of target vectors.
\end{proof}

\paragraph{What is left.}
The nearconstant contribution has a separate direct stub-exposure
proof in the sibling research note
\texttt{parity-reduction/nearconstant-configuration-atoms.tex}.
Combining that result with Theorem~\ref{thm:tails} leaves only the
central arbitrary-output local limit theorem, including its lattice
index, feasibility condition, leading constant and uniform error.
Neither the zero-target theorem of Huang nor the present entropy
bounds alone supplies that final step. No odd-block RSBM simplicity
theorem is claimed in this note.

\begin{thebibliography}{9}
\bibitem{huang} J. Huang,
\emph{Invertibility of adjacency matrices for random $d$-regular graphs},
arXiv:1807.06465v2, December 30, 2018.
\url{https://arxiv.org/pdf/1807.06465v2}.
The paper appeared in Duke Mathematical Journal 170(18), 3977--4032
(2021); the precise version inspected here is the cited preprint.
\end{thebibliography}
\end{document}
